\begin{ex}
Gieo một con xúc xắc 6 mặt (từ 1 đến 6 chấm). Có thể xuất hiện mặt 7 chấm không? (Có ghi 1, Không ghi 0).
\shortans{0}
\loigiai{
Xúc xắc chỉ có từ 1 đến 6 chấm nên không thể có 7 chấm.
}
\end{ex}

\begin{ex}
Gieo một đồng xu có hai mặt Sấp (S) và Ngửa (N). Có bao nhiêu khả năng có thể xảy ra về mặt xuất hiện?
\shortans{2}
\loigiai{
Có 2 khả năng: mặt Sấp hoặc mặt Ngửa.
}
\end{ex}

\begin{ex}
Trong hộp chỉ có 5 viên bi đỏ. Lấy ngẫu nhiên 1 viên bi thì chắc chắn lấy được bi màu gì? (Bi đỏ ghi 1, Bi xanh ghi 2).
\shortans{1}
\loigiai{
Chắc chắn lấy được bi đỏ.
}
\end{ex}

\begin{ex}
Trong hộp có 3 viên bi đỏ và 3 viên bi xanh. Lấy ngẫu nhiên 1 viên bi, có thể lấy được viên bi vàng không? (Có ghi 1, Không thể ghi 0).
\shortans{0}
\loigiai{
Không thể lấy được bi vàng.
}
\end{ex}

\begin{ex}
Trong hộp có 4 viên bi xanh và 2 viên bi đỏ. Sự kiện 'Lấy được viên bi xanh' là: (Chắc chắn ghi 1, Có thể ghi 2, Không thể ghi 3).
\shortans{2}
\loigiai{
Có thể lấy được bi xanh (hoặc bi đỏ).
}
\end{ex}

\begin{ex}
Mặt trời mọc ở hướng Đông vào mỗi sáng là sự kiện gì? (Chắc chắn ghi 1, Có thể ghi 2, Không thể ghi 3).
\shortans{1}
\loigiai{
Chắc chắn xảy ra.
}
\end{ex}

\begin{ex}
Một con cá biết bay trên trời như chim là sự kiện gì? (Chắc chắn ghi 1, Có thể ghi 2, Không thể ghi 3).
\shortans{3}
\loigiai{
Không thể xảy ra.
}
\end{ex}

\begin{ex}
Ngày mai trời có mưa là sự kiện gì? (Chắc chắn ghi 1, Có thể ghi 2, Không thể ghi 3).
\shortans{2}
\loigiai{
Có thể xảy ra.
}
\end{ex}

\begin{ex}
Một hộp có 10 chiếc kẹo, trong đó có 10 chiếc kẹo cam. Lấy ngẫu nhiên 1 chiếc kẹo thì chắc chắn lấy được kẹo gì? (Kẹo cam ghi 1, kẹo dâu ghi 2).
\shortans{1}
\loigiai{
Chắc chắn là kẹo cam.
}
\end{ex}

\begin{ex}
Tung một con xúc xắc, sự kiện 'Số chấm xuất hiện nhỏ hơn 7' là sự kiện gì? (Chắc chắn ghi 1, Có thể ghi 2, Không thể ghi 3).
\shortans{1}
\loigiai{
Chắc chắn, vì xúc xắc có từ 1 đến 6 chấm (đều < 7).
}
\end{ex}

\begin{ex}
Tung một con xúc xắc, sự kiện 'Số chấm xuất hiện bằng 0' là sự kiện gì? (Chắc chắn ghi 1, Có thể ghi 2, Không thể ghi 3).
\shortans{3}
\loigiai{
Không thể, xúc xắc không có mặt 0 chấm.
}
\end{ex}

\begin{ex}
Tung một con xúc xắc, sự kiện 'Số chấm xuất hiện là số chẵn' là sự kiện gì? (Chắc chắn ghi 1, Có thể ghi 2, Không thể ghi 3).
\shortans{2}
\loigiai{
Có thể xảy ra (mặt 2, 4, hoặc 6).
}
\end{ex}

\begin{ex}
Có bao nhiêu mặt là số chẵn trên một con xúc xắc 6 mặt?
\shortans{3}
\loigiai{
Có 3 mặt chẵn: 2, 4, 6.
}
\end{ex}

\begin{ex}
Có bao nhiêu mặt là số lẻ trên một con xúc xắc 6 mặt?
\shortans{3}
\loigiai{
Có 3 mặt lẻ: 1, 3, 5.
}
\end{ex}

\begin{ex}
Trong túi có 5 quả bóng xanh và 5 quả bóng đỏ. Không nhìn vào túi, phải lấy ra ít nhất bao nhiêu quả bóng để chắc chắn có ít nhất 1 quả bóng đỏ?
\shortans{6}
\loigiai{
Trường hợp xấu nhất lấy phải 5 quả bóng xanh, cần lấy quả thứ 6 để chắc chắn có bóng đỏ.
}
\end{ex}

\begin{ex}
Trong hộp có 4 viên bi vàng và 6 viên bi đỏ. Cần lấy ra ít nhất bao nhiêu viên bi để chắc chắn có ít nhất 1 viên bi vàng?
\shortans{7}
\loigiai{
Trường hợp xấu nhất lấy hết 6 viên bi đỏ, lấy thêm 1 viên nữa: $6 + 1 = 7$ viên.
}
\end{ex}

\begin{ex}
Trong hộp có 10 viên bi xanh. Lấy ra 2 viên bi thì có thể có viên bi đỏ nào không? (Có ghi 1, Không ghi 0).
\shortans{0}
\loigiai{
Không thể có bi đỏ.
}
\end{ex}

\begin{ex}
Quay một vòng quay chia thành 4 phần bằng nhau có ghi số 1, 2, 3, 4. Có bao nhiêu khả năng mũi tên dừng lại ở một ô số?
\shortans{4}
\loigiai{
Có 4 khả năng.
}
\end{ex}

\begin{ex}
Vòng quay có 4 số: 1, 2, 3, 4. Mũi tên có thể chỉ vào số 5 không? (Có ghi 1, Không ghi 0).
\shortans{0}
\loigiai{
Không thể chỉ vào số 5.
}
\end{ex}

\begin{ex}
Rút ngẫu nhiên một tấm thẻ từ tập thẻ đánh số từ 1 đến 10. Có thể rút được thẻ ghi số 8 không? (Có ghi 1, Không ghi 0).
\shortans{1}
\loigiai{
Có thể rút được thẻ số 8.
}
\end{ex}

\begin{ex}
Rút ngẫu nhiên một thẻ từ tập thẻ từ 1 đến 10. Sự kiện rút được thẻ số 12 là sự kiện gì? (Chắc chắn ghi 1, Có thể ghi 2, Không thể ghi 3).
\shortans{3}
\loigiai{
Không thể xảy ra.
}
\end{ex}

\begin{ex}
Rút ngẫu nhiên một thẻ từ 1 đến 10. Có bao nhiêu thẻ mang số chẵn?
\shortans{5}
\loigiai{
Các số 2, 4, 6, 8, 10 (5 thẻ).
}
\end{ex}

\begin{ex}
Rút ngẫu nhiên một thẻ từ 1 đến 10. Có bao nhiêu thẻ mang số lẻ?
\shortans{5}
\loigiai{
Các số 1, 3, 5, 7, 9 (5 thẻ).
}
\end{ex}

\begin{ex}
Bạn Nam tung đồng xu 10 lần thì số lần mặt Ngửa xuất hiện có thể lớn hơn 10 không? (Có ghi 1, Không ghi 0).
\shortans{0}
\loigiai{
Tung 10 lần thì tối đa chỉ có 10 lần ngửa.
}
\end{ex}

\begin{ex}
Hai bạn chơi oẳn tù tì. Trong một ván chơi có mấy kết quả có thể xảy ra đối với bạn thứ nhất? (Thắng, thua, hòa).
\shortans{3}
\loigiai{
Có 3 kết quả: Thắng, Thua hoặc Hòa.
}
\end{ex}
